\section*{Acknowledgements}
This paper makes use of the following ALMA data: the \NProjCodes{} projects
listed in Table~\ref{tab:projcodes}, each cited as \texttt{ADS/JAO.ALMA\#}
followed by its project code. The \NEB{} primary-census execution blocks
searched span \PrEbDateFirst{} to \PrEbDateLast{}, and we thank the principal
investigators and teams who proposed and obtained
them.\ProjCodeUnresClause{} ALMA is a partnership of ESO
(representing its member states), NSF (USA) and NINS (Japan), together with NRC
(Canada), NSTC and ASIAA (Taiwan), and KASI (Republic of Korea), in cooperation
with the Republic of Chile. The Joint ALMA Observatory is operated by ESO,
AUI/NRAO and NAOJ. This work has made use of data from the European Space
Agency (ESA) mission {\it Gaia} \citep{GaiaDR3}, processed by the {\it Gaia}
Data Processing and Analysis Consortium (DPAC). This research has made use
of the SIMBAD database \citep{SIMBAD2000} (CDS, Strasbourg, France); of the
NASA Exoplanet Archive \citep{NEA2013}, which is operated by the California
Institute of Technology under contract with the National Aeronautics and
Space Administration under the Exoplanet Exploration Program; and of the
Splatalogue spectral-line database \citep{Splatalogue2007}.

\emph{Use of generative artificial intelligence.} Large language model
assistance (Anthropic Claude) was used in manuscript preparation of this
complex paper. Data analysis was done using the Python based framework ASTRA
\citep[][\href{https://arxiv.org/abs/2610.36144}{https://\allowbreak arxiv.org/\allowbreak abs/\allowbreak 2610.36144}]{White2026ASTRA}
to assist in the archival retrieval and calibration scripts, search and
vetting pipeline, statistical tests, and the generators that produce every
number, and carefully cross-checked and validated. The scientific content,
analysis, and conclusions were designed, executed, and verified by the
authors. All errors and omissions are the responsibility of the author.

\emph{Testing and validation.} No result here is accepted because a model
reported it: every number in the text is written by the program that computes it from the
data, and a separate audit re-derives the identities that must hold among the
published values and refuses the paper where one does not. Every check in the
analysis was first made to fail on deliberately wrong input before it was
relied on.

\emph{Software.} The pipeline depends on \textsc{casa}~6.7.0
\citep{CASA2022}, \textsc{python}~3 with \textsc{astropy} \citep{Astropy2022}, \textsc{numpy}
\citep{Harris2020} and \textsc{scipy} \citep{Virtanen2020}, and the ALMA Science
Archive and SIMBAD TAP interfaces through \textsc{astroquery} \citep{Ginsburg2019};
all are free and open-source, at the versions cited.

\emph{Author contributions.} G.J.W.\ conceived the survey, defined the target
sample and the selection criteria, and wrote the manuscript; R.D.\ built and
operated the archival retrieval, calibration and search pipeline, implemented
the spatial control statistic and the injection--recovery campaigns, and
produced the per-window catalogue and its regeneration scripts. Both authors
designed the vetting procedure, chose the masked species, carried out the
statistical validation reported in \S\ref{sec:statistic} and
Appendix~\ref{app:falsealarm}, and agreed the dispositions of the crossing
ledger.
%% Authors to confirm this division before submission.

\emph{Funding.} This work received no dedicated funding. It used public archival data and
no observing time was awarded for it.
%% Authors to replace this if any grant supported the work.

\emph{Competing interests.} The authors declare no competing interests. VBRL Holdings is a privately
held company through which R.D.\ conducts independent research, named
here solely as his affiliation; it had no role in the design of the
survey, the selection of targets, the analysis, the decision to publish
or the content of this paper, and provided no funding for the work.
%% Authors to confirm this wording, which is accurate only if the company's
%% involvement was indeed nil.

\section*{Data availability}
\label{sec:data}
Every observation used here is public in the ALMA Science Archive under the
\NProjCodes{} project codes of Table~\ref{tab:projcodes}, and is retrieved by
the execution-block identifier, which the crossing ledger gives for each of the
\EvNCrossRaw{} threshold crossings \citep{Chen2022}; the astrometry is
Gaia~DR3, the planet counts come from the NASA Exoplanet Archive, and the line
mask from one Splatalogue query of CDMS and JPL made on \McQueryDate. The data
and software are held in a versioned Zenodo record \citep{SIdeposit}. It
contains the per-window catalogue of all \NWindows{} windows, the Class~A and
Class~B limits with the injection trials behind them, the crossing ledger, the
control-field statistics, the execution-block and session accounting (every
row of Table~\ref{tab:blockledger}), and the
resampling inputs and seeds. From these one program,
\texttt{code/reproduce.py}, rebuilds the chance expectation and its
probability, the completeness factor, the sensitivity summaries and the
repeat-window comparison, and checks each against the value printed here; the
principal tables also accompany the arXiv version as ancillary files. The
record holds the search pipeline as it ran, which takes a public execution
block to the statistics of this paper. Some windows, in blocks that were not
recalibrated, carry products made by the forms of the extractor and of the
search that preceded the final ones. The final extractor places the phase
centre at the star's parallax-displaced position and keeps every field entry,
which the first forms did not, and the very first also drew its control
annulus differently. Those versions are archived with their settings,
together with a table that gives, for every window, the versions that made
its products; run again on freshly calibrated 7\,m and 12\,m blocks, they
reproduce those products in every window tested. The dynamic spectra at the
star and at its \RcNRetCtrl{} retained control positions, around each of the
\DaNSpecCross{} crossings, form a second file of the same record
(\DaSpecFileMB\,MB). The calibrated visibilities, of order a terabyte, and the
per-window spectra were not kept, for lack of storage, so this execution of
the search can be run again from the archive with the same code but cannot be
repeated exactly.
